\begin{frame}{계산 이어가기: 스케일링·softmax·가중합}
  $\sqrt{d_k} = \sqrt{2}$ 로 나눈 뒤 각 행에 softmax:
  \begin{itemize}
    \item 행 1: softmax(0.354, 0) = \textbf{(0.587, 0.413)}
      \; (0.354 $= 0.5/\sqrt{2}$)
    \item 행 2: softmax(0, 0.707) = \textbf{(0.330, 0.670)}
      \; (0.707 $= 1/\sqrt{2}$)
  \end{itemize}
  \vskip0.5em
  최종 출력 ($V_1=(2,1), V_2=(1,2)$):
  \begin{itemize}
    \item 단어 1: $0.587\,V_1 + 0.413\,V_2 = \mathbf{(1.587,\ 1.413)}$
    \item 단어 2: $0.330\,V_1 + 0.670\,V_2 = \mathbf{(1.330,\ 1.670)}$
  \end{itemize}
  \vskip0.5em
  \begin{block}{두 가지 관찰}
    (1) \textbf{단어 1의 출력이 $V_1$ 쪽에, 단어 2의 출력이 $V_2$
    쪽에 기울어 있다} (0.59, 0.67) — ``관련 단어의 내용을 관련도만큼
    섞는다''가 숫자로 성립. (2) 출력이 순수한 복사
    (정답이면 (2,1)·(1,2))는 \textbf{아니다} — 다른 단어의 내용이
    조금 섞여 있다.
  \end{block}
  \vskip0.4em
  어텐션 출력은 ``가장 관련 있는 단어 \textbf{하나}를 고르는
  것''이 아니라 \textbf{``모든 Value의 가중평균''} — 이 체감이
  있으면 12.2절의 softmax가 ``\textbf{왜 합이 1인 분포}여야 하는
  지''가 자동으로 이해된다.
  \vskip0.3em
  \textbf{점수 차이가 클수록 어텐션 가중치는 더 집중}된다:
  단어 2의 점수 간격(1 vs 0)이 단어 1(0.5 vs 0)보다 크므로
  가중치도 더 기울어 있다 (0.67/0.33 vs 0.59/0.41). 어떤 Key의
  점수가 압도적으로 높으면 softmax는 가중치의 거의 전부를 거기에
  몰아넣는다 — 12.2절 실습의 ``그것은''이 ``동물''에 0.56을
  주는 현상은 바로 이 메커니즘의 \textbf{더 극단적인 버전}.
\end{frame}
